% AUTHOR: Amlal El Mahrouss
% PURPOSE: WG04.2: Sum Analysis.

\documentclass[11pt, a4paper]{article}
\usepackage{graphicx}
\usepackage{listings}
\usepackage{amsmath,amssymb,amsthm}
\usepackage{xcolor}
\usepackage{hyperref}
\usepackage{mathtools}
\usepackage[margin=0.5in,top=1in,bottom=1in]{geometry}

\title{Sum Analysis for Integrals.}
\author{Amlal El Mahrouss\\\texttt{amlal@nekernel.org}}
\date{February 2026}

\begin{document}

\maketitle

\begin{center}
	\rule[0.01cm]{17cm}{0.01cm}
\end{center}

\begin{abstract}
    The paper covers several definitions that are made to be used for sum analysis; their purpose has been kept abstract for flexibility reasons.
\end{abstract}

\section{DEFINITIONS}

We assume $x \in \mathbb{R}, \quad a, b \in \mathbb{N}, \quad b > a > 0$ for:

\subsection{DEFINITION OF THE S-INTEGRAL:}
\begin{equation}
    I_{1}(a, b, u, n) \coloneqq \int_{a}^{b} S(t, u, n)dt
\end{equation}
\subsection{DEFINITION OF THE S-FUNCTION:}
\begin{equation}
    S(t, u, n) = \sum_{k=1}^{n} (\frac{dt}{du}\cdot u), \quad t > 0, \quad n > k
\end{equation}

\section{INTRODUCTION}

We will develop our conditions before going to the second part of the paper.

\end{document}

